% =============================================================================
%  SBTTD Mach Post-Processing GUI — User & Developer Manual
%  Author: Hadie Hesham Sabbah
%  Last updated: May 2026
% =============================================================================
\documentclass[11pt,letterpaper]{article}

% ---- packages ----
\usepackage[margin=1in]{geometry}
\usepackage{graphicx}
\usepackage{xcolor}
\usepackage{hyperref}
\usepackage{listings}
\usepackage{booktabs}
\usepackage{amsmath,amssymb}
\usepackage{enumitem}
\usepackage{fancyhdr}
\usepackage{titlesec}
\usepackage[T1]{fontenc}
\usepackage{lmodern}
\usepackage{parskip}

% ---- colours ----
\definecolor{codebg}{RGB}{245,245,245}
\definecolor{codeframe}{RGB}{200,200,200}
\definecolor{keyword}{RGB}{0,100,180}
\definecolor{string}{RGB}{180,40,40}
\definecolor{comment}{RGB}{100,140,100}
\definecolor{sectionblue}{RGB}{30,70,130}

% ---- listings style ----
\lstset{
  language=Python,
  basicstyle=\ttfamily\small,
  backgroundcolor=\color{codebg},
  frame=single,
  rulecolor=\color{codeframe},
  keywordstyle=\color{keyword}\bfseries,
  stringstyle=\color{string},
  commentstyle=\color{comment}\itshape,
  breaklines=true,
  tabsize=4,
  showstringspaces=false,
  numbers=left,
  numberstyle=\tiny\color{gray},
  numbersep=8pt,
}

% ---- section formatting ----
\titleformat{\section}{\Large\bfseries\color{sectionblue}}{}{0em}{}
\titleformat{\subsection}{\large\bfseries\color{sectionblue!80}}{}{0em}{}
\titleformat{\subsubsection}{\normalsize\bfseries\color{sectionblue!60}}{}{0em}{}

% ---- header / footer ----
\pagestyle{fancy}
\fancyhf{}
\fancyhead[L]{\small SBTTD GUI Manual}
\fancyhead[R]{\small\thepage}
\renewcommand{\headrulewidth}{0.4pt}

% ---- title ----
\title{%
  \vspace{-1cm}
  \textbf{SBTTD Mach Post-Processing GUI}\\[6pt]
  \large User \& Developer Manual
}
\author{Hadie Hesham Sabbah}
\date{May 2026}

\begin{document}
\maketitle
\tableofcontents
\newpage

% =============================================================================
\section{Introduction}
% =============================================================================

The \textbf{SBTTD Mach Post-Processing GUI} provides a graphical interface for
the entire CFD post-processing pipeline used in the Supersonic Bladeless Turbine
Theoretical Development (SBTTD) project.  It replaces the need to step through
individual code cells in \texttt{MachPostProcessCode.py} by organising the
workflow into discrete, interactive tabs.

The GUI is built with \textbf{customtkinter} (a modern-looking wrapper around
Tkinter) and embeds \textbf{Matplotlib} figures directly inside the application
window for interactive plotting.

\subsection{What the GUI Does}

\begin{enumerate}[leftmargin=2em]
  \item Load raw Tecplot binary files (\texttt{.bin}) or previously saved
        pickle data.
  \item Extract wall variables (pressure, shear stress, Mach number, etc.)
        with configurable spatial masks.
  \item Compute boundary-layer thickness via Tecplot line extraction.
  \item Detect flow separation, compute separation lengths, and generate
        heatmaps.
  \item Plot any wall variable against streamwise position, filtered by
        \(h/l\) or Mach number.
  \item Compute axial force, torque, and power from 2-D pressure integration.
  \item Run theoretical models (small-perturbation theory, shock-expansion
        theory) and compare against CFD.
  \item Export Mach contour images and visualise total-pressure-loss maps.
\end{enumerate}

% =============================================================================
\section{Installation \& Dependencies}
% =============================================================================

\subsection{Python Environment}

The GUI requires Python~3.9 or later.  The recommended setup is a dedicated
conda or virtualenv environment.

\begin{lstlisting}[language=bash]
# Create environment (conda example)
conda create -n sbttd python=3.11 -y
conda activate sbttd

# Install dependencies
pip install customtkinter matplotlib numpy scipy xarray pandas
pip install pygasflow sympy seaborn
pip install loguru python-dotenv typer

# Tecplot (requires a licence)
pip install pytecplot
\end{lstlisting}

\subsection{Required Packages}

\begin{table}[h]
\centering
\begin{tabular}{lll}
\toprule
\textbf{Package} & \textbf{Purpose} & \textbf{Required?} \\
\midrule
customtkinter    & GUI framework                 & Yes \\
matplotlib       & Plotting                      & Yes \\
numpy            & Numerical computation         & Yes \\
scipy            & Interpolation, root finding   & Yes \\
xarray           & Structured datasets           & Yes \\
pandas           & DataFrames (force tables)     & Yes \\
pytecplot        & Tecplot integration           & For Tabs 1, 3, 8 \\
pygasflow        & Isentropic solvers            & For Tab 7 \\
sympy            & Symbolic math                 & For Tab 7 \\
\bottomrule
\end{tabular}
\caption{Dependency matrix.}
\end{table}

\subsection{Launching the GUI}

Navigate to the \texttt{2\_SBTTD Data Analysis Code} directory and run:

\begin{lstlisting}[language=bash]
python -m gui
\end{lstlisting}

This invokes \texttt{gui/\_\_main\_\_.py}, which creates the main
\texttt{SBTTDApp} window and starts the Tkinter event loop.

% =============================================================================
\section{Architecture Overview}
% =============================================================================

\subsection{Directory Layout}

\begin{lstlisting}[language=bash,numbers=none]
2_SBTTD Data Analysis Code/
|-- gui/
|   |-- __init__.py
|   |-- __main__.py          # Entry point
|   |-- app.py               # Main window (SBTTDApp)
|   |-- state.py             # Central application state (AppState)
|   |-- tabs/
|   |   |-- __init__.py
|   |   |-- tab_data_io.py       # Tab 1: Data Import / Export
|   |   |-- tab_variables.py     # Tab 2: Variable Extraction
|   |   |-- tab_boundary_layer.py# Tab 3: Boundary Layer Analysis
|   |   |-- tab_separation.py    # Tab 4: Separation Analysis
|   |   |-- tab_wall_plots.py    # Tab 5: Wall Variable Plots
|   |   |-- tab_force_power.py   # Tab 6: Force & Power
|   |   |-- tab_theoretical.py   # Tab 7: Theoretical Models
|   |   |-- tab_contours.py      # Tab 8: Contour Export
|   |-- widgets/
|       |-- __init__.py
|       |-- plot_canvas.py   # Embedded Matplotlib canvas
|       |-- dir_browser.py   # Directory/file browser widget
|-- utils/
|   |-- dataload_util.py     # bigImport, runLoader, runSaver
|   |-- parameterComputation.py # variableImporterMasked, ReCompute
|   |-- models.py            # Shock-expansion, force, separation
|   |-- plotting.py          # Publication-quality plotting functions
|   |-- experimental.py      # Experimental data processing
|   |-- computation.py       # (placeholder)
|-- MachPostProcessCode.py   # Original script (reference)
|-- sweep_study_gui.py       # Standalone exp. comparison GUI
\end{lstlisting}

\subsection{Design Principles}

\begin{itemize}[leftmargin=2em]
  \item \textbf{Central state object.}  All data lives in a single
        \texttt{AppState} instance (\texttt{gui/state.py}).  Every tab reads
        from and writes to this object, eliminating the need to pass data
        between tabs explicitly.

  \item \textbf{Tab-per-workflow-step.}  Each tab corresponds to one logical
        step in the post-processing pipeline, matching the cell groups in the
        original \texttt{MachPostProcessCode.py}.

  \item \textbf{Threaded long operations.}  Data import, variable extraction,
        and boundary-layer computation run in background threads so the GUI
        remains responsive.

  \item \textbf{Embedded Matplotlib.}  The \texttt{PlotCanvas} widget wraps
        \texttt{FigureCanvasTkAgg} with a navigation toolbar, allowing
        interactive pan/zoom/save directly inside the app.
\end{itemize}

\subsection{The AppState Class}

\texttt{gui/state.py} defines \texttt{AppState} with the following attribute
groups:

\begin{table}[h]
\centering
\small
\begin{tabular}{ll}
\toprule
\textbf{Group}         & \textbf{Key Attributes} \\
\midrule
Raw data               & \texttt{ds\_by\_case}, \texttt{ds\_by\_case\_quad}, \texttt{ds\_by\_case\_inlet} \\
Wall variables         & \texttt{x}, \texttt{y}, \texttt{P}, \texttt{tau\_x}, \texttt{tau\_y}, \texttt{mach}, \ldots \\
Quadrant variables     & \texttt{x\_quad}, \texttt{P\_quad}, \texttt{mach\_quad}, \ldots \\
Inlet variables        & \texttt{P\_inlet}, \texttt{mach\_inlet}, \texttt{u\_inlet} \\
Computed quantities    & \texttt{tau\_wall}, \texttt{delta\_n\_dict} \\
Separation results     & \texttt{sep\_length}, \texttt{x\_sep}, \texttt{x\_attach}, \ldots \\
Force / theory results & \texttt{df\_axial\_force}, \texttt{axialForceScaled}, \ldots \\
Case metadata          & \texttt{case\_keys}, \texttt{h\_l\_values}, \texttt{mach\_values}, \texttt{cases\_by\_hl} \\
Status flags           & \texttt{data\_loaded}, \texttt{variables\_extracted}, \texttt{tecplot\_connected} \\
\bottomrule
\end{tabular}
\caption{AppState attribute summary.}
\end{table}

% =============================================================================
\section{Tab-by-Tab Usage Guide}
% =============================================================================

\subsection{Tab 1: Data I/O}

\textbf{Purpose:} Import raw Tecplot files or load previously saved pickle
data; optionally save the current dataset.

\textbf{Workflow:}
\begin{enumerate}
  \item Click \emph{Connect to Tecplot} (required for raw import and BL analysis).
  \item Either:
    \begin{itemize}
      \item Browse to a study root directory and click \emph{Run bigImport}, or
      \item Browse to a folder containing \texttt{.pkl} files and click
            \emph{Load Data}.
    \end{itemize}
  \item After loading, the status label shows the number of cases found.
  \item Use \emph{Save Data} to serialise the current datasets to pickle files.
\end{enumerate}

\textbf{Backend function:} \texttt{utils.dataload\_util.bigImport},
\texttt{runLoader}, \texttt{runSaver}.

% ----------------------------------------------------------------------------
\subsection{Tab 2: Variable Extraction}

\textbf{Purpose:} Extract wall variables from \texttt{ds\_by\_case} with an
optional spatial mask, and compute $\tau_{\text{wall}}$.

\textbf{Parameters:}
\begin{itemize}
  \item \texttt{min\_l}, \texttt{max\_l}: streamwise bounds for the mask
        (set both to 0 to disable masking).
\end{itemize}

\textbf{Output:} Populates \texttt{state.x}, \texttt{state.P},
\texttt{state.tau\_x}, etc.

\textbf{Backend function:} \texttt{utils.parameterComputation.variableImporterMasked}.

% ----------------------------------------------------------------------------
\subsection{Tab 3: Boundary Layer Analysis}

\textbf{Purpose:} Compute BL thickness by extracting line probes normal to the
wall in Tecplot and detecting the vorticity-gradient edge.

\textbf{Requires:} Tecplot connection, extracted variables.

\textbf{Parameters:}
\begin{itemize}
  \item \texttt{Ny}: number of points per probe (default 1000).
  \item \texttt{BL height}: maximum probe length in mm (default 3).
  \item \texttt{Stride}: skip every $n$-th wall point (default 1).
  \item \texttt{Threshold}: normalised vorticity-gradient threshold (default 0.02).
\end{itemize}

\textbf{Output:} \texttt{state.delta\_n\_dict} — dictionary of BL thickness
arrays in mm.

% ----------------------------------------------------------------------------
\subsection{Tab 4: Separation Analysis}

\textbf{Purpose:} Identify separation and reattachment locations from the sign
of $\tau_{\text{wall}}$, compute total separation length, and generate summary
plots.

\textbf{Steps:}
\begin{enumerate}
  \item Click \emph{Compute Separation} $\rightarrow$ calls
        \texttt{find\_sepLength}.
  \item Optionally click \emph{Find Geometry Extrema} to locate wave peaks
        and troughs.
  \item Select a plot type and click \emph{Plot}:
    \begin{itemize}
      \item \textbf{Separation Heatmap} — binary matrix (sep/no-sep) coloured
            green/red.
      \item \textbf{Sep Length vs h/l} — dimensional separation length curves.
      \item \textbf{Sep Length (non-dim) vs h/l} — $L_{\text{sep}}/L_{\text{domain}}$.
      \item \textbf{BL \& Separation Contours} — scatter of separation
            percentage in $(h/l, M)$ space.
    \end{itemize}
\end{enumerate}

\textbf{Backend function:} \texttt{utils.models.find\_sepLength},
\texttt{max\_min\_finder}.

% ----------------------------------------------------------------------------
\subsection{Tab 5: Wall Variable Plots}

\textbf{Purpose:} Quick plotting of any extracted wall variable ($\tau_w$,
$y^+$, $P$, $\partial P/\partial x$, Mach, etc.) against $x$.

\textbf{Controls:}
\begin{itemize}
  \item \textbf{Variable}: dropdown of available quantities.
  \item \textbf{Filter by}: \texttt{h\_l} or \texttt{mach} $+$ a numeric
        value.  Cases matching that parameter are overlaid; the varying
        parameter is used for legend labels.
  \item \textbf{Plot}: single-axes plot for the selected filter value.
  \item \textbf{Subplot All h/l}: one subplot per $h/l$, all Mach curves in
        each.
\end{itemize}

% ----------------------------------------------------------------------------
\subsection{Tab 6: Force \& Power}

\textbf{Purpose:} Integrate wall pressure to obtain axial force, then compute
torque and power for given rotor parameters.

\textbf{Parameters:}
\begin{itemize}
  \item Rotor radius $R$ [m], RPM, span [m].
\end{itemize}

\textbf{Steps:}
\begin{enumerate}
  \item Click \emph{Compute Force (2D)}.
  \item Click \emph{Compute Power} (uses $P = F \cdot R \cdot \omega$).
  \item Optionally click \emph{Build Force DataFrame} for a tabular summary.
  \item Select a plot type and click \emph{Plot}.
\end{enumerate}

\textbf{Backend:} \texttt{utils.models.compute\_force\_2D},
\texttt{create\_axial\_force\_dataframe}.

% ----------------------------------------------------------------------------
\subsection{Tab 7: Theoretical Models (Shock-Expansion)}

\textbf{Purpose:} Run small-perturbation theory and/or shock-expansion theory
to predict pressure distributions and axial force, and compare to CFD.

\textbf{Models:}
\begin{itemize}
  \item \textbf{Small Perturbation}: linearised supersonic thin-airfoil theory.
  \item \textbf{Shock-Expansion (SE)}: marching oblique-shock / Prandtl-Meyer
        expansion analysis.
  \item \textbf{Combined}: hybrid SP $+$ SE model.
\end{itemize}

\textbf{Plots:}
\begin{itemize}
  \item Scaled force ratio $F_{\text{RANS}}/F_{\text{theory}}$ vs $h/l$.
  \item First-shock pressure vs Mach.
  \item Pressure distribution overlay.
\end{itemize}

\textbf{Backend:} \texttt{utils.models.smallPertSolver},
\texttt{smallPertSolver\_with\_SE}, \texttt{get\_first\_shock\_pressures}.

% ----------------------------------------------------------------------------
\subsection{Tab 8: Contour Export}

\textbf{Purpose:} Export Mach contour images from Tecplot and visualise
total-pressure-loss maps.

\textbf{Sections:}
\begin{itemize}
  \item \textbf{Mach Contour Export}: provide study directory, PNG output
        directory, and optional Tecplot layout file; click \emph{Export
        Contours}.
  \item \textbf{Total Pressure Loss}: scatter map of pressure loss in
        $(h/l, M)$ space, computed from quad/inlet data.
  \item \textbf{Mach Contour Visualisation}: assemble exported PNGs into a
        subplot grid for a given $h/l$.
\end{itemize}

\textbf{Backend:} \texttt{utils.plotting.export\_mach\_contours},
\texttt{plot\_mach\_contours\_per\_hl}.

% =============================================================================
\section{How to Modify \& Extend the GUI}
% =============================================================================

\subsection{Adding a New Tab}

\begin{enumerate}
  \item Create a new file \texttt{gui/tabs/tab\_myfeature.py}.
  \item Define a class inheriting from \texttt{ctk.CTkFrame}:
\begin{lstlisting}
import customtkinter as ctk

class MyFeatureTab(ctk.CTkFrame):
    def __init__(self, master, state, status_callback=None, **kwargs):
        super().__init__(master, **kwargs)
        self.state = state
        self.status = status_callback or (lambda msg: None)
        self._build_ui()

    def _build_ui(self):
        # Create your widgets here
        ctk.CTkLabel(self, text="My Feature").pack()
        ctk.CTkButton(self, text="Run", command=self._run).pack()

    def _run(self):
        # Access shared data via self.state
        if self.state.x is None:
            return
        # ... your logic ...
        self.status("Done!")
\end{lstlisting}

  \item Register the tab in \texttt{gui/app.py}:
\begin{lstlisting}
from gui.tabs.tab_myfeature import MyFeatureTab

# inside SBTTDApp._build_tabs(), add to tab_specs:
("9. My Feature", MyFeatureTab),
\end{lstlisting}
\end{enumerate}

\subsection{Adding a New State Variable}

Open \texttt{gui/state.py} and add the attribute in
\texttt{AppState.\_\_init\_\_}:

\begin{lstlisting}
self.my_new_variable = None
\end{lstlisting}

All tabs automatically have access via \texttt{self.state.my\_new\_variable}.

\subsection{Adding a New Widget}

Place reusable widgets in \texttt{gui/widgets/}.  The existing
\texttt{PlotCanvas} and \texttt{DirBrowser} serve as templates:

\begin{itemize}
  \item Inherit from \texttt{ctk.CTkFrame}.
  \item Accept \texttt{master} as the first argument.
  \item Provide getter/setter methods for the widget's value.
\end{itemize}

\subsection{Adding a New Backend Function}

Place computation logic in the appropriate \texttt{utils/} module:

\begin{itemize}
  \item \texttt{utils/parameterComputation.py} — variable extraction, Reynolds
        number.
  \item \texttt{utils/models.py} — theoretical models, separation, force.
  \item \texttt{utils/plotting.py} — publication-quality plotting.
  \item \texttt{utils/dataload\_util.py} — data I/O, Tecplot import.
  \item \texttt{utils/experimental.py} — experimental data processing.
\end{itemize}

Then call the function from your tab, preferably inside a
\texttt{threading.Thread} for long-running operations to keep the GUI
responsive.

\subsection{Changing the Theme}

The sidebar includes a theme toggle (Dark / Light / System).
Programmatically, call:

\begin{lstlisting}
ctk.set_appearance_mode("Dark")   # or "Light", "System"
ctk.set_default_color_theme("blue")  # or "green", "dark-blue"
\end{lstlisting}

% =============================================================================
\section{Backend Utility Reference}
% =============================================================================

\subsection{\texttt{utils/dataload\_util.py}}

\begin{table}[h]
\centering
\small
\begin{tabular}{lp{8cm}}
\toprule
\textbf{Function} & \textbf{Description} \\
\midrule
\texttt{bigImport(base\_dir, fileName)} &
  Recursively load Tecplot \texttt{.bin} files from a study directory.
  Returns \texttt{(ds\_by\_case, ds\_by\_case\_quad, ds\_by\_case\_inlet)}.\\
\texttt{runSaver(ds, ds\_quad, ds\_inlet)} &
  Pickle the three dictionaries to a date-stamped folder.\\
\texttt{runLoader(load\_dir\_dic)} &
  Load the latest dated pickle folder. Returns the three dictionaries.\\
\texttt{file\_pathFinder(fileName, rootDir)} &
  Locate all instances of \texttt{fileName} under \texttt{rootDir}.\\
\texttt{load\_minfo\_step\_force(path, use)} &
  Parse mcfd info files for convergence data.\\
\bottomrule
\end{tabular}
\end{table}

\subsection{\texttt{utils/parameterComputation.py}}

\begin{table}[h]
\centering
\small
\begin{tabular}{lp{8cm}}
\toprule
\textbf{Function} & \textbf{Description} \\
\midrule
\texttt{variableImporterMasked(ds, min\_l, max\_l, mask\_input)} &
  Extract 18 wall variables from \texttt{ds\_by\_case}, optionally masked to
  $[\texttt{min\_l}, \texttt{max\_l}]$.\\
\texttt{ReCompute(T, U, rho, \ldots)} &
  Compute Reynolds number per case using Sutherland viscosity.\\
\texttt{yplusThreshold(y\_plus)} &
  Warn if any case exceeds $y^+ > 1$.\\
\bottomrule
\end{tabular}
\end{table}

\subsection{\texttt{utils/models.py} (selected)}

\begin{table}[h]
\centering
\small
\begin{tabular}{lp{8cm}}
\toprule
\textbf{Function} & \textbf{Description} \\
\midrule
\texttt{find\_sepLength(ds, x, tau\_wall)} &
  Detect separation/reattachment from $\tau_w$ sign changes; returns
  lengths, locations.\\
\texttt{compute\_force\_2D(x, y, P, span)} &
  Integrate pressure to get axial force per unit span.\\
\texttt{smallPertSolver(\ldots)} &
  Small-perturbation theory for supersonic wavy walls.\\
\texttt{get\_first\_shock\_pressures(\ldots)} &
  Shock-expansion analysis returning the pressure at the first shock.\\
\texttt{max\_min\_finder(ds, x, y)} &
  Locate geometry maxima and minima (wave peaks/troughs).\\
\bottomrule
\end{tabular}
\end{table}

\subsection{\texttt{utils/plotting.py} (selected)}

\begin{table}[h]
\centering
\small
\begin{tabular}{lp{8cm}}
\toprule
\textbf{Function} & \textbf{Description} \\
\midrule
\texttt{plotter(x, y, \ldots)} &
  Single-curve publication plot.\\
\texttt{plotter\_multiPerCase(\ldots)} &
  Multi-curve plot filtered by $h/l$ or Mach.\\
\texttt{subplotter\_multiPerCase(\ldots)} &
  Subplot grid, one per filter value.\\
\texttt{plot\_BL\_thickness(\ldots)} &
  BL thickness curves, one figure per $h/l$.\\
\texttt{export\_mach\_contours(\ldots)} &
  Batch-export Mach contour images from Tecplot.\\
\bottomrule
\end{tabular}
\end{table}

% =============================================================================
\section{Typical Workflow}
% =============================================================================

A recommended session follows these steps:

\begin{enumerate}
  \item \textbf{Tab 1} — Connect to Tecplot (if needed); load data via
        \emph{bigImport} or \emph{Load Data}.
  \item \textbf{Tab 2} — Set \texttt{min\_l = 0}, \texttt{max\_l = 0.1};
        click \emph{Extract Variables}, then \emph{Compute tau\_wall}.
  \item \textbf{Tab 3} — Set BL parameters; click \emph{Compute BL Thickness}
        (requires Tecplot); plot results.
  \item \textbf{Tab 4} — Click \emph{Compute Separation}; view heatmap and
        length plots.
  \item \textbf{Tab 5} — Browse wall variables ($\tau_w$, $y^+$, $P$, etc.)
        interactively.
  \item \textbf{Tab 6} — Compute axial force and power.
  \item \textbf{Tab 7} — Run small-perturbation or SE model; compare to CFD.
  \item \textbf{Tab 8} — Export contours; view pressure-loss maps.
\end{enumerate}

% =============================================================================
\section{Troubleshooting}
% =============================================================================

\begin{description}[leftmargin=2em]
  \item[``Tecplot connection failed'']
    Ensure Tecplot 360 is running and PyTecplot is installed.  The GUI calls
    \texttt{tp.session.connect()}, which requires the Tecplot TecUtil server.

  \item[GUI freezes during long computation]
    All heavy functions run in daemon threads.  If the GUI appears
    unresponsive, wait for the status bar to update.  Do not close the window
    during computation.

  \item[ModuleNotFoundError for \texttt{customtkinter}]
    Run \texttt{pip install customtkinter}.

  \item[Plots appear blank]
    Make sure data is loaded (Tab~1) and variables are extracted (Tab~2)
    before attempting to plot in later tabs.

  \item[``No cases match h\_l = \ldots'']
    Case keys follow the pattern \texttt{h\_l\_0.04\_Mach\_2.0}.  Ensure
    the filter value matches the numeric format in your dataset.
\end{description}

% =============================================================================
\section{Appendix: Key Equations}
% =============================================================================

\subsection{Wall Shear Stress}

The tangential wall shear stress is computed from the Cartesian components:
\begin{equation}
  \tau_w = \tau_x \, t_x + \tau_y \, t_y
\end{equation}
where $\mathbf{t} = (t_x, t_y) = (\mathrm{d}x/\mathrm{d}s,\;
\mathrm{d}y/\mathrm{d}s)$ is the unit tangent vector along the wall.

\subsection{Boundary Layer Edge Detection}

The BL edge is located where the normalised vorticity gradient drops below a
threshold $\epsilon$:
\begin{equation}
  \frac{|\partial \omega_z / \partial n|}{|\partial \omega_z / \partial n|_{\max}}
  < \epsilon, \quad \epsilon = 0.02 \text{ (default)}
\end{equation}

\subsection{Separation Length}

Separation is detected where $\tau_w < 0$.  The total separation length is:
\begin{equation}
  L_{\text{sep}} = \sum_{i} (x_{\text{attach},i} - x_{\text{sep},i})
\end{equation}
The non-dimensional form is $L_{\text{sep}} / L_{\text{domain}}$.

\subsection{Axial Force (2-D Pressure Integration)}

\begin{equation}
  F_{\text{axial}} = -\int_0^L p(x) \frac{\mathrm{d}y}{\mathrm{d}x}\,\mathrm{d}x
  \;\times\; b
\end{equation}
where $b$ is the rotor span (out-of-plane depth).

\subsection{Prandtl--Meyer Function}

\begin{equation}
  \nu(M) = \sqrt{\frac{\gamma+1}{\gamma-1}}\;
  \arctan\!\sqrt{\frac{\gamma-1}{\gamma+1}(M^2-1)}
  \;-\; \arctan\!\sqrt{M^2-1}
\end{equation}

\end{document}
